\documentclass[10pt,a4paper]{amsart}
\usepackage[margin=19mm]{geometry}
\usepackage{amsmath,amssymb,mathtools}
\usepackage[scr=rsfs,cal=euler]{mathalfa}
\usepackage{xfrac}
\usepackage[unicode,hidelinks,bookmarks=false]{hyperref}
\newcommand{\ie}{i.e.\ }
\newcommand{\cf}{cf.\ }
\newcommand{\resp}{respectively\ }
\newcommand{\Subsection}{\subsection}
\providecommand{\nobreakdash}{\nobreak\hbox{-}}
\setlength{\emergencystretch}{2em}
\multlinegap0pt
\usepackage{bm}
\usepackage{bbm}
\usepackage{dsfont}
\usepackage{xspace}
\usepackage{cancel}
\usepackage{mathtools}
\usepackage{amsthm}
\makeatletter
\@ifundefined{prop}{\newtheorem{prop}{Proposition}}{}
\@ifundefined{thm}{\newtheorem{thm}{Theorem}}{}
\makeatother
\newcommand{\inner}[2]{\left( #1, #2 \right)}
\title[Convergent Operator-Splitting Scheme for Viscosity Solutions]{Convergent Operator-Splitting Scheme for Viscosity Solutions: A Foundation for Learning Domain-to-Solution Maps}
\author{Po-Yi Wu}
\date{Source: arXiv:2505.20618v5 (CC BY 4.0)}
\begin{document}
\maketitle

\noindent\emph{Source and licence.} Convergent Operator-Splitting Scheme for Viscosity Solutions: A Foundation for Learning Domain-to-Solution Maps. Version arXiv:2505.20618v5 (CC BY 4.0), distributed under the Creative Commons Attribution 4.0 International licence, as recorded in the arXiv licence metadata for this exact version and shown on the abstract page https://arxiv.org/abs/2505.20618v5. Full rights notice: licenses/arxiv-2505.20618-CC-BY-4.0.txt.\par\smallskip

\noindent\textit{Editorial scope.} Counted unit: the Thm whose source label is \texttt{thm:convergence\_rate\_narrative} in the article (source0.tex lines 716--722, source label \texttt{thm:convergence\_rate\_narrative}), delivered together with the article's own complete original proof of it (source0.tex lines 724--740, one \texttt{proof} environment of 17 lines). No mathematics is written, repaired or re-derived: statement and proof are the article's own text line for line. Required context retained uncounted: eq:general\_system (lines 207--214); eq:predictor\_step (lines 423--426); eq:corrector\_step (lines 430--436); prop:strong\_truncation\_error (lines 519--524); thm:dcp\_proved (lines 603--607). Every definition, display and label of those lines is retained unchanged. Omitted from this package and not redistributed: 1--206, 215--422, 427--429, 437--518, 525--602, 608--715, 723--723, 741--1069. Nothing inside a retained excerpt is omitted or altered: every retained body is delivered exactly as the article's lines are written, so no omission receipt of any kind accompanies this package. Citations: the retained excerpts carry no citation command at all, so no bibliography entry of the article is transcribed and no reference is redistributed. Reference adaptation: none was needed and none was made; every reference inside the delivered unit resolves inside it, so no unresolved reference or citation is produced. Printed slips kept literal: no printed word, symbol, hypothesis or number of the counted statement or of its proof was altered. Layout and class: the article's own class is \texttt{sn-jnl} with packages including graphicx, multirow, amsmath, amssymb, amsfonts, amsthm, mathrsfs, appendix, xcolor, textcomp, manyfoot, booktabs, algorithm2e, algorithmicx; the wrapper is the lane's pinned \texttt{amsart} editorial wrapper at 10pt on A4, which restates the theorem-like environments the retained text needs and adds the article's own macro definitions that the retained text needs (\texttt{\textbackslash inner}), with extra wrapper packages (bm, bbm, dsfont, xspace, cancel, mathtools). The delivered numbering is the unit's own. Assets: no retained span contains a figure environment, a graphics-inclusion command, a PDF-page inclusion, a table or a TikZ picture; no image file is copied into this package, the package declares no asset dependency and no image file is redistributed with it. Licence: version arXiv:2505.20618v5 of this article is distributed under the Creative Commons Attribution 4.0 International licence, as recorded in the arXiv licence metadata for this exact version and shown on the abstract page https://arxiv.org/abs/2505.20618v5. A full rights notice is delivered at licenses/arxiv-2505.20618-CC-BY-4.0.txt. Characters: every retained excerpt is pure ASCII, so the delivered engine, XeLaTeX with its default Latin Modern text font, reports no missing character.
\begin{equation}
\label{eq:general_system}
\begin{cases}
\partial_t y = f(\bm{x}, t, y, \nabla y, \nabla^2 y) + \bm{B}(\bm{x}, t, \nabla y) \cdot \bm{p} & \text{in } \Omega_T, \\
\bm{g}(\bm{x}, t, y, \nabla y) = 0 & \text{in } \Omega_T, \\
y = y_0 & \text{on } \partial_p \Omega_T,
\end{cases}
\end{equation}

    \begin{equation}
    \label{eq:predictor_step}
        \left\langle \frac{\hat{y}_h^{m+1} - y_h^m}{\Delta t}, v_h \right\rangle = \left\langle f_h(\bm{x}, t_m, y_h^m, \nabla y_h^m), v_h \right\rangle, \quad \forall v_h \in V_h.
    \end{equation}

    \begin{equation}
    \label{eq:corrector_step}
    \begin{cases}
        \left\langle \frac{y_h^{m+1} - \hat{y}_h^{m+1}}{\Delta t}, v_h \right\rangle + \mathcal{A}_h(y_h^{m+1}, v_h) = \left\langle \bm{B}(\bm{x}, t_{m+1}, \nabla y_h^{m+1}) \cdot \bm{p}_h^{m+1}, v_h \right\rangle, \\
        \left\langle \bm{g}(\bm{x}, t_{m+1}, y_h^{m+1}, \nabla y_h^{m+1}), \bm{q}_h \right\rangle = 0.
    \end{cases}
    \end{equation}

\begin{prop}[Local Truncation Error Bound]
\label{prop:strong_truncation_error}
Let the exact solution $y$ of \eqref{eq:general_system} possess higher regularity, specifically $y \in C([0,T]; H^{3}(\Omega)) \cap C^1([0,T]; H^1(\Omega)) \cap C^2([0,T]; L^2(\Omega))$. The local truncation error $\tau_h^{m+1}$, defined by substituting the exact solution $y$ into the scheme, is bounded in the dual norm of $H^1(\Omega)$ as:
$$ \| \tau_h^{m+1} \|_{(H^1(\Omega))'} \leq C (\Delta t + h^2), $$
where the constant $C$ depends on norms of the exact solution but is independent of $h$ and $\Delta t$.
\end{prop}

\begin{thm}[Discrete Comparison Principle]
\label{thm:dcp_proved}
Let $\{y_h^m\}$ and $\{z_h^m\}$ be two solutions of the scheme defined by \eqref{eq:predictor_step} and the linearized version of \eqref{eq:corrector_step}, corresponding to initial and boundary data that are ordered such that $y_h^0 \le z_h^0$ and $y_h^m|_{\partial\Omega} \le z_h^m|_{\partial\Omega}$ for all $m \ge 0$. Suppose that the operator $f(x, t, y, \bm{q}, \bm{M})$ is non-increasing with respect to its state variable $y$ with Lipschitz constant $L_f$, that the time-derivative is discretized using a \textbf{lumped mass matrix} $M_h$, that the stabilization parameter $\sigma$ satisfies the condition derived in the proof, and that the time step $\Delta t$ satisfies the condition $\Delta t \le \min_i (M_{ii}) / L_f$. Then the scheme satisfies the discrete comparison principle,
$$y_h^m \le z_h^m \quad \text{for all } m \ge 0.$$
\end{thm}

\begin{thm}[Convergence Rate]
\label{thm:convergence_rate_narrative}
Let the conditions of the Discrete Comparison Principle (Theorem \ref{thm:dcp_proved}) hold. If the unique viscosity solution possesses additional regularity, specifically $y \in C^1([0,T]; C^3(\overline{\Omega})) \cap C^2([0,T]; C^2(\overline{\Omega}))$, then the numerical scheme converges with a rate governed by the discretization parameters. For a time step $\Delta t$ and mesh size $h$ satisfying $\Delta t \le Ch$, there exists a constant $C'$, independent of $h$ and $\Delta t$, such that the error is bounded as
\[
\max_{0 \le m \le T/\Delta t} \| y(\cdot, t_m) - y_h^m \|_{L^2(\Omega)} \leq C' (\Delta t + h^2).
\]
\end{thm}

\begin{proof}
The analysis hinges on bounding the departure of the numerical solution $y_h^m$ from the exact solution $y^m = y(\cdot, t_m)$. We decompose the total error into an approximation error $\eta^m$ from the projection and a discrete error $\theta_h^m$ using the elliptic projection $\mathcal{P}_h y^m \in V_h$ of the exact solution:
$$y^m - y_h^m = (y^m - \mathcal{P}_h y^m) - (y_h^m - \mathcal{P}_h y^m) = \eta^m - \theta_h^m.$$
Standard finite element approximation theory provides the bound $\|\eta^m\|_{L^2(\Omega)} \le C h^3$ for $P_2$ elements, given the assumed smoothness of $y$. The core of the proof is to control the evolution of the discrete error $\theta_h^m$.

We derive an evolution equation for $\theta_h^m$ by subtracting the numerical scheme's equation for $y_h^{m+1}$ from the same discrete relation satisfied by the projected exact solution, $\mathcal{P}_h y^{m+1}$. This process isolates the local truncation error, $\tau_h^{m+1}$, which represents the extent to which the exact solution fails to satisfy the scheme. The resulting error equation takes the form
$$\inner{\frac{\theta_h^{m+1} - \theta_h^m}{\Delta t}}{v_h} + \mathcal{A}_h(\theta_h^{m+1}, v_h) = \langle \tau_h^{m+1}, v_h \rangle + \langle \mathcal{N}_h, v_h \rangle,$$
where $\mathcal{N}_h$ encapsulates nonlinear error terms. Under the stated regularity assumptions, we can directly apply Proposition~\ref{prop:strong_truncation_error} to bound the local truncation error:
$$ \| \tau_h^{m+1} \|_{(H^1(\Omega))'} \le C(\Delta t + h^2). $$
The established stability of the scheme is the critical tool to control the accumulation of these local errors. A standard stability analysis, obtained by choosing $v_h = \theta_h^{m+1}$ as the test function, exploits the coercivity of the form $\mathcal{A}_h$ and the $L^\infty$-bound on the solution to control the nonlinearities. This procedure yields a recursive inequality for the discrete error:
$$\|\theta_h^{m+1}\|^2_{L^2} \le (1 + C_1 \Delta t) \|\theta_h^m\|^2_{L^2} + C_2 \Delta t (\Delta t + h^2)^2.$$
An application of the discrete Gronwall's lemma, assuming the initial error $\theta_h^0 = 0$, reveals that the accumulated discrete error is controlled over $[0,T]$, yielding $\|\theta_h^m\|_{L^2(\Omega)} \le C'_T (\Delta t + h^2)$.

Finally, the total error is bounded by the triangle inequality:
$$\|y^m - y_h^m\|_{L^2(\Omega)} \le \|\eta^m\|_{L^2(\Omega)} + \|\theta_h^m\|_{L^2(\Omega)} \le C h^3 + C'_T(\Delta t + h^2).$$
As the $O(\Delta t + h^2)$ term is dominant, the desired convergence rate is established.
\end{proof}

\end{document}
